\section{Introduction}

The spatial spread of infectious diseases depends both on pathogen biology and on the structure of human mobility. Characteristics such as transmissibility, infectious period, latency, mortality, reinfection, or waning immunity determine how an outbreak evolves locally, while transport infrastructures determine how infections move between populations. Air, sea, and land transport systems have increased the reach, speed, and volume of human movement, making mobility networks central to the geographical dissemination of infectious diseases \cite{tatem2006global}. In highly connected regions such as Europe, local outbreaks can therefore become regional or continental, with air travel supporting rapid long-distance seeding \cite{colizza:airline} and multiscale mobility processes, including commuting and long-range travel, shaping spatial and temporal spreading patterns \cite{balcan:multiscale}.

Mathematical epidemic models provide a controlled way to compare these mechanisms. Classical compartmental models divide the population into epidemiological states, beginning with the SIR framework of Kermack and McKendrick \cite{kermack:sir}, in which susceptible individuals become infected and are then removed or recovered. Extensions represent different disease archetypes: SEIR adds a latent exposed state, SIS represents infections without lasting immunity, SEIRS allows immunity to wane, and SEIQRD includes quarantine and disease-induced death \cite{paul2023fractional}. Such models are widely used to analyse transmission mechanisms, epidemic thresholds, and control strategies \cite{hethcote2000mathematics}. However, homogeneous compartmental models describe only the dynamics within a single population. To model spatial spread between cities, this work uses a metapopulation approach, where each city is a local population with internal epidemic dynamics and mobility flows connect cities, allowing infections to be transported across the network \cite{colizza2008metapopulation}.

Figure~\ref{fig:metapopulation_framework} summarizes this modelling logic by combining within-city epidemic dynamics with mobility flows between cities.

\begin{figure*}[tp]
    \centering
    \includegraphics[width=0.66\textwidth]{figures/metapopulation_epidemic_network_diagram.png}
    \caption[Metapopulation epidemic framework]{
    Metapopulation epidemic framework. Each circle represents a city, modelled as a local subpopulation with its own compartmental epidemic dynamics. The internal compartments illustrate a generic SEIR-type progression: $S$ denotes susceptible individuals, $E$ exposed individuals who are infected but not yet infectious, $I$ infectious individuals, and $R$ recovered or removed individuals. The dashed return arrow indicates that some models, such as SIS or SEIRS, allow individuals to become susceptible again. Arrows between cities represent mobility flows that can transport susceptible, exposed, or infected individuals and seed outbreaks elsewhere. Inspired by metapopulation epidemic models such as Colizza and Vespignani \cite{colizza2008metapopulation}.
    }
    \label{fig:metapopulation_framework}
\end{figure*}

Europe is represented here as a multilayer transport network composed of air, water, and land mobility layers. Multilayer network theory is appropriate for systems in which the same set of nodes can be connected by different types of links, each corresponding to a distinct interaction channel or transport mode \cite{kivela2014multilayer,boccaletti2014structure}. In this work, cities are nodes and transport routes are organised by layer: air links represent fast long-distance mobility, water links represent maritime and ferry-based connectivity, and land links represent regional travel and commuting. This distinction is epidemiologically relevant because mobility can operate through different mechanisms. Air and water travel are approximated as more diffusive, meaning that individuals may relocate and mix in a destination city, whereas land commuting is recurrent, meaning that individuals travel, mix temporarily, and return to their home city. Recurrent mobility can synchronize nearby cities, while long-distance diffusive mobility can seed infection in distant populations \cite{belik2011mobility,balcan:multiscale}.

Figure~\ref{fig:diffusive_recurrent_mobility} illustrates the conceptual difference between these two mobility mechanisms.

\begin{figure}[tbp]
    \centering
    \includegraphics[width=0.82\columnwidth]{figures/mobility_mechanisms_diffusive_vs_recurrent_types.png}
    \caption[Diffusive and recurrent mobility mechanisms]{
    Diffusive and recurrent mobility mechanisms. In diffusive mobility, individuals move from city $i$ to city $j$, relocate, and mix with the destination population; this approximates air travel and some water-based connections. In recurrent mobility, individuals travel from a home city $i$ to a destination $j$, mix temporarily, and return home; this represents land commuting and regional travel. Diffusive mobility can seed distant outbreaks, whereas recurrent mobility mainly couples and synchronizes nearby cities \cite{belik2011mobility,balcan:multiscale}.
    }
    \label{fig:diffusive_recurrent_mobility}
\end{figure}

Beyond mobility volume, the topology of the network influences epidemic propagation. Cities and routes are not structurally equivalent: some cities act as hubs with many connections, others as bridges between regions, and others belong to dense structural cores that may sustain spreading processes. Epidemic processes in complex networks are strongly affected by heterogeneity, connectivity, and node position \cite{pastorsatorras:review}. Moreover, geographical distance is not always the best proxy for epidemiological proximity, since strongly connected cities can be effectively close even when they are geographically distant \cite{brockmann2013hidden}.

Network science metrics can therefore support intervention design. Degree identifies the highest-traffic cities, betweenness centrality identifies the cities that bridge regions, and dense-neighbourhood measures such as triangle participation identify cities embedded in tightly connected cores that can sustain spreading \cite{kitsak:spreaders,milo:motifs}. More recent optimal-percolation methods, notably Collective Influence and non-backtracking centrality, go further by targeting the low-degree, high-impact bridge cities that raw degree overlooks \cite{morone:influence,radicchi:nonbacktracking}. These metrics can guide node interventions, interpreted as vaccination, protection, screening, or intensified surveillance in selected cities, and edge interventions, interpreted as closure or restriction of selected routes. Targeted immunization and edge-removal strategies can outperform random interventions in heterogeneous networks, although their effectiveness depends on the network structure and the timing and intensity of control \cite{pastorsatorras:immunization,bajardi2011mobility,yang2013edge}.

The objective of this work is comparative rather than predictive: the aim is not to reproduce a specific European epidemic or to provide operational forecasts, but to compare how different compartmental models, mobility layers, and intervention strategies shape simulated spreading. The central question is how the multilayer structure of European transport enables or limits different epidemic processes, and how network science metrics can guide effective containment. Combining compartmental models, a metapopulation framework, multilayer networks, and topology-based interventions, the study systematically relates human mobility, network structure, and the spatial dynamics of infectious disease spread.